\section{Classification of real branch points and maximal-modulus roles}\label{sec:ep-classification}

Among the five finite branch values identified in Theorem~\ref{thm:elliptic-reduction}, exactly three are real.
By Lemma~\ref{lem:cyclic-companion}, every repeated root of \(\Pi(\cdot,y)\) corresponds to a defective eigenvalue of the companion family \(T(y)\).
In this section we determine their spectral role by analyzing the eigenvalue collision structure at each branch point and, crucially, by comparing the modulus of the colliding eigenvalues against the remaining roots.

The following elementary comparison is the only general local fact needed below:
a branch that leaves a repeated root at order \(t^{1/\nu}\) dominates any non-bifurcating simple branch whose modulus changes only at order \(t\).

\begin{lemma}[Local dominance from a fixed square-root splitting]\label{lem:local-dominance}
Let \(P(\lambda,t)=\prod_{j=1}^n(\lambda-\mu_j(t))\) be a monic degree-\(n\) polynomial whose coefficients are analytic in \(t\) near \(t=0\).
The conclusion is local on a prescribed real side of the base: throughout the statement take \(t>0\) sufficiently small and use the positive real square-root coordinate \(s=t^{1/2}>0\).
Suppose that at \(t=0\) exactly two roots collide:
\[
\mu_1(0)=\mu_2(0)=:\mu_*,
\qquad
\mu_3(0),\dots,\mu_n(0) \text{ simple and distinct from } \mu_*.
\]
Assume that for all \(j\ge 3\),
\begin{equation}\label{eq:local-dominance-hyp}
|\mu_j(0)|\le |\mu_*|.
\end{equation}
Assume that, after this square-root base change, one of the two bifurcating roots is represented by a holomorphic branch \(\widetilde\mu_1(s)\), and write \(\mu_1(t)=\widetilde\mu_1(t^{1/2})\) for the chosen positive-side branch.
Assume that this branch satisfies
\begin{equation}\label{eq:bifurc-form}
\abs{\mu_1(t)} \ge \abs{\mu_*} + a t^{1/2} - B t
\end{equation}
for some constants \(a,B>0\) and all real \(t>0\) sufficiently small on that side.
Then there exists \(\delta>0\) such that
\begin{equation}\label{eq:local-dominance-concl}
\abs{\widetilde\mu_1(s)}>\abs{\mu_j(s^2)}
\qquad\text{for every } j\ge 3 \text{ and } 0<s<\delta^{1/2}.
\end{equation}
\end{lemma}

\begin{proof}
Since the non-colliding roots \(\mu_3(0),\dots,\mu_n(0)\) are simple, the implicit function theorem gives analytic continuations \(\mu_j(t)\) on a two-sided neighborhood of \(0\) and constants \(L_j>0\) such that
\[
\abs{\mu_j(t)-\mu_j(0)}\le L_j t
\qquad (j\ge 3,\ 0\le t\le t_0)
\]
for some \(t_0>0\).
\[
\abs{\mu_j(t)}
\le
\abs{\mu_j(0)}+L_j t
\le
\abs{\mu_*}+L_j t
=
\abs{\mu_*}+O(t)
\qquad (j\ge 3).
\]
On the other hand, the chosen square-root branch satisfies, with \(s=t^{1/2}>0\),
\[
\abs{\widetilde\mu_1(s)}=\abs{\mu_1(t)}\ge \abs{\mu_*}+a t^{1/2}-Bt.
\]
Let \(L:=\max_{j\ge 3}L_j\).
Since \(t^{1/2}\gg t\) as \(t\to0^+\), choose \(0<\delta\le t_0\) so that
\[
Bt\le \frac{a}{2}t^{1/2},
\qquad
Lt\le \frac{a}{2}t^{1/2}
\qquad (0<t<\delta).
\]
Then for \(0<t<\delta\),
\[
\abs{\widetilde\mu_1(s)}
\ge
\abs{\mu_*}+a t^{1/2}-Bt
\ge
\abs{\mu_*}+\frac{a}{2}t^{1/2}
>
\abs{\mu_*}+Lt
\ge
\abs{\mu_j(t)}
\]
for every \(j\ge 3\), proving \eqref{eq:local-dominance-concl} after writing \(t=s^2\).
\end{proof}

\begin{theorem}[Canonical specialization of the fixed-member real-projection criterion (Theorem B)]\label{thm:theorem-b-real-projection-package}
For the canonical quartic
\[
\Pi(\lambda,y)=\lambda^4-\lambda^3-(2y+1)\lambda^2+\lambda+y(y+1)
\]
and its elliptic normalization
\[
E:\quad Y^2=4\lambda^3-4\lambda+1,
\qquad
y=\lambda^2-\frac12+\frac12Y,
\]
the real-projection package is the following single chain of assertions.
\begin{enumerate}
\item The fixed-member finite criterion is Theorem~\ref{thm:balanced-real-projection}: for each real component of one smooth balanced member, the real \(y\)-image and real-fibre transition skeleton are determined by the finite algebraic set \(\mathcal V(J)\) of endpoint values, non-endpoint ramification values, and the exceptional \(\lambda=0\) values.
\item In the specialization \((a,c,d,g)=(-1,-1,1,0)\), Proposition~\ref{prop:quartic-real-projection-skeleton} gives the complete obstruction skeleton
\[
\mathcal V(J_{\mathrm c})=
\left\{\beta_1^2-\frac12,\ \beta_2^2-\frac12,\ 1,\ y_{\mathrm{sub}}\right\},
\qquad
\mathcal V(J_{\mathrm n})=
\left\{\beta_3^2-\frac12,\ 0\right\},
\]
with extrema
\[
\min \mathcal V(J_{\mathrm c})=y_{\mathrm{sub}},
\qquad
\max \mathcal V(J_{\mathrm c})=1,
\qquad
\min \mathcal V(J_{\mathrm n})=0.
\]
Hence
\[
y(E_{\mathrm c}(\RR))=[y_{\mathrm{sub}},1],
\qquad
y(E_{\mathrm n}(\RR))=[0,\infty).
\]
\item For every regular real value of \(y\), the source of each real root is fixed by the table
\[
\begin{array}{c|c|c|c}
\text{range of }y & \#E_{\mathrm c} & \#E_{\mathrm n} & \#\Pi_y(\RR)\\
\hline
y<y_{\mathrm{sub}} & 0 & 0 & 0\\
y_{\mathrm{sub}}<y<0 & 2 & 0 & 2\\
0<y<1 & 2 & 2 & 4\\
y>1 & 0 & 2 & 2.
\end{array}
\]
The component decomposition behind this count is recorded in Corollary~\ref{cor:real-phase-diagram}.
\end{enumerate}
No recurrence-theoretic or dominance input enters this package; those inputs begin only with Theorem~\ref{thm:consolidated-branch-dominance-package}.
\end{theorem}

\begin{proof}
The first assertion is precisely Theorem~\ref{thm:balanced-real-projection}, applied before specialization.
For \((a,c,d,g)=(-1,-1,1,0)\), Proposition~\ref{prop:quartic-real-projection-skeleton} computes the two real intervals of \(4\lambda^3-4\lambda+1\ge0\), lists all endpoint and critical values in the two finite sets above, and orders those algebraic values so that the compact and noncompact oval images are respectively \([y_{\mathrm{sub}},1]\) and \([0,\infty)\).
Let now \(y\) be a regular real value.
On a real oval, the projection is a continuous degree-two real covering over the interior of its image: away from the displayed branch values no point is ramified, so an interior value contributes two real points from that oval and an exterior value contributes none.
Applying this to the two intervals \([y_{\mathrm{sub}},1]\) and \([0,\infty)\) gives the four rows of the table.
The equality between real points on the normalization and real roots of \(\Pi(\cdot,y)\) follows from the normalization map away from branch values, where all roots are simple.
Corollary~\ref{cor:real-phase-diagram} retains the component source of the calculation, and the stated dependency chain is complete.
\end{proof}

\begin{lemma}[Real-root count for the residual branch cubic]\label{lem:cubic-branch-real-root}
The cubic
\[
c(y)=256y^3+411y^2+165y+32
\]
has exactly one real root.
This root lies in \((-2,-1)\), and more sharply in
\begin{equation}\label{eq:ysub-certified-interval}
-\frac{2269}{2000}<y<-\frac{5672}{5000}.
\end{equation}
\end{lemma}

\begin{proof}
For a cubic, the sign of the discriminant determines the number of real roots: a negative discriminant gives one real root and one non-real conjugate pair.
For
\[
c(y)=ay^3+by^2+cy+d,
\qquad
(a,b,c,d)=(256,411,165,32),
\]
the cubic discriminant
\[
\Delta=b^2c^2-4ac^3-4b^3d-27a^2d^2+18abcd
\]
is \(\Delta=-699868431<0\).
Thus \(c\) has exactly one real zero.
Finally,
\[
c(-2)=-702<0,
\qquad
c(-1)=22>0,
\]
so the intermediate value theorem places that unique real zero in \((-2,-1)\).
The refined interval follows from the rational evaluations
\[
c\!\left(-\frac{2269}{2000}\right)=-\frac{5301369}{500000000}<0,
\qquad
c\!\left(-\frac{5672}{5000}\right)=\frac{2804026}{244140625}>0,
\]
and uniqueness of the real zero.
\end{proof}

\begin{definition}[The subdominant real branch value]\label{def:ysub}
Throughout the paper, \(\ySub\) denotes the unique real root of
\begin{equation}\label{eq:ysub-definition}
256\ySub^3+411\ySub^2+165\ySub+32=0,
\qquad
-\frac{2269}{2000}<\ySub<-\frac{5672}{5000}.
\end{equation}
Equivalently, \(\ySub\in(-1.1345,-1.1344)\).  This is the negative real branch value that later occurs in the real oval images and in the real-spectral phase diagram.
\end{definition}

\begin{lemma}[Exact subdominance certificate at \(\ySub\)]\label{lem:ysub-exact-subdominance}
Let \(\alpha\) be the unique real root of
\[
h(\lambda)=16\lambda^3-9\lambda^2+1.
\]
Then \(-11/40<\alpha<-27/100\), and the real branch value \(\ySub\) is
\begin{equation}\label{eq:ysub-alpha-formula}
\ySub=\frac{4\alpha^3-3\alpha^2-2\alpha+1}{4\alpha}.
\end{equation}
At \(y=\ySub\) the spectral polynomial has the exact factorization in \(\QQ(\alpha)[\lambda]\)
\begin{equation}\label{eq:ysub-factorization}
\Pi(\lambda,\ySub)
=(\lambda-\alpha)^2
\left(
\lambda^2+\frac{(\alpha-1)(\alpha+1)}{8\alpha^2}\lambda
+\frac{9\alpha^2-144\alpha-1}{256\alpha^2}
\right).
\end{equation}
The quadratic factor is irreducible over \(\RR\), and if its two roots are \(z,\overline z\), then
\begin{equation}\label{eq:ysub-modulus-gap}
|z|^2
=\frac{9\alpha^2-144\alpha-1}{256\alpha^2}
>\alpha^2.
\end{equation}
Consequently the colliding double root at \(y=\ySub\) is strictly subdominant.
\end{lemma}

\begin{proof}
The derivative \(h'(\lambda)=3\lambda(16\lambda-6)\) has critical points \(0\) and \(3/8\).
Since \(h(0)=1\), \(h(3/8)=37/64>0\), and \(h(\lambda)\to-\infty\) as \(\lambda\to-\infty\), the cubic \(h\) has exactly one real root.
The rational evaluations
\[
h\!\left(-\frac{11}{40}\right)=-\frac{107}{8000}<0,
\qquad
h\!\left(-\frac{27}{100}\right)=\frac{7243}{250000}>0
\]
therefore isolate it in \((-11/40,-27/100)\).
The critical-value formula from the elliptic projection gives \eqref{eq:ysub-alpha-formula}, and substituting \eqref{eq:ysub-alpha-formula} into \(c(y)\) reduces its numerator to zero modulo \(h(\alpha)\); since \(\alpha\ne0\), this identifies the value with the unique real root \(\ySub\) of \(c\).

A direct division in the cubic field gives \eqref{eq:ysub-factorization}.
Equivalently, after multiplying by \(16\alpha^2\), the quotient of \(\Pi(\lambda,\ySub)\) by \((\lambda-\alpha)^2\) has coefficients
\[
16\alpha^2,
\qquad
2(\alpha-1)(\alpha+1),
\qquad
\frac{9\alpha^2-144\alpha-1}{16},
\]
and the remainder is zero in \(\QQ[\alpha]/(h)\).
The discriminant of the displayed quadratic is
\[
\frac{(\alpha-1)^2(\alpha+1)^2}{64\alpha^4}
-\frac{9\alpha^2-144\alpha-1}{64\alpha^2}
=\frac{2479\alpha^2+16\alpha-247}{8(81\alpha^2-16\alpha-9)}.
\]
On the isolating interval \((-11/40,-27/100)\), the numerator is negative and the denominator is positive; indeed
\[
2479\left(-\frac{11}{40}\right)^2+16\left(-\frac{11}{40}\right)-247=-\frac{102281}{1600}<0,
\]
while
\[
81\left(-\frac{27}{100}\right)^2-16\left(-\frac{27}{100}\right)-9=\frac{12249}{10000}>0.
\]
The derivative of the numerator, \(4958\lambda+16\), is negative throughout the interval, so the numerator stays negative after its left-endpoint check.  The denominator is the upward-opening quadratic \(81\lambda^2-16\lambda-9\), and its positive zero is larger than \(1/3\); hence the denominator is positive throughout \((-11/40,-27/100)\).
Thus the quadratic has negative discriminant and its two roots are a non-real conjugate pair.
Their squared modulus is the constant term of the quadratic, and
\[
|z|^2-\alpha^2
=\frac{9\alpha^2-144\alpha-1}{256\alpha^2}-\alpha^2
=-\frac{9\alpha^2+16\alpha-1}{32\alpha^2}.
\]
The numerator \(9\lambda^2+16\lambda-1\) has derivative \(18\lambda+16>0\) on \((-11/40,-27/100)\), and its value at \(-27/100\) is \(-46639/10000<0\); hence it is negative at \(\alpha\).
Since \(32\alpha^2>0\), the last display is positive.
Therefore the two non-real simple roots have modulus strictly larger than \(|\alpha|\), proving strict subdominance of the double root.
\end{proof}

\begin{theorem}[Real branch-point classification]\label{thm:ep-classification}
The three real branch values are
\[
\ySub < 0, \qquad y_0 = 0, \qquad y_1 = 1,
\]
where \(\ySub\) is the algebraic number fixed in Definition~\ref{def:ysub}.
Their spectral roles are as follows.
\begin{enumerate}
\item At $y = 0$, the spectral polynomial factors as
\[
\Pi(\lambda, 0) = \lambda(\lambda + 1)(\lambda - 1)^2.
\]
The double root $\lambda=1$ ties in modulus with the simple root $\lambda=-1$.
Nevertheless, for $y>0$ sufficiently small, the pair of branches issuing from $\lambda=1$ contains the unique maximal-modulus root.
Thus $y=0$ is the endpoint of the maximal-modulus branch over $y>0$.

\item At $y = 1$, the spectral polynomial factors as
\[
\Pi(\lambda, 1) = (\lambda - 2)(\lambda - 1)(\lambda + 1)^2.
\]
The double root $\lambda = -1$ is not the maximal-modulus eigenvalue; the positive real root $\lambda = 2$ remains simple.
Hence $y = 1$ is not a branch point of the maximal-modulus branch over $y>0$.

\item At \(y=\ySub\), the colliding double root is \emph{strictly subdominant}: it is exceeded in modulus by a pair of complex-conjugate roots.
Therefore $y_{\mathrm{sub}}$ is a subdominant branch point, not on the maximal-modulus envelope.
\end{enumerate}
\end{theorem}

\begin{proof}
Lemma~\ref{lem:cubic-branch-real-root} shows that the residual branch cubic \(c(y)\) contributes exactly one real branch value, and that it is negative.
Together with the two real branch values \(0\) and \(1\) from Theorem~\ref{thm:elliptic-reduction}, this gives precisely the three real branch values displayed in the theorem.

We begin with the local dominance claim in part~(1).

\medskip
\noindent\textbf{Continuity of roots at $y=0$.}
At $y=0$ the factorization is $\Pi(\lambda,0)=\lambda(\lambda+1)(\lambda-1)^2$.
The three simple roots $\lambda\in\{0,-1\}$ extend analytically in $y$ by the implicit function theorem, since $\partial_\lambda\Pi$ is nonzero at each of them.
The double root at $\lambda=1$ produces a square-root ramification as described next.

\medskip
\noindent\textbf{Self-contained local expansion.}
Write $y=s^2$ for real $s>0$ small.
Set $\lambda=1+h$ directly in $\Pi(\lambda,y)=\lambda^4-\lambda^3-(2y+1)\lambda^2+\lambda+y(y+1)$ and collect terms:
\begin{align*}
\Pi(1+h,y)
&= (1+4h+6h^2+4h^3+h^4) - (1+3h+3h^2+h^3)\\
&\quad - (1+2y)(1+2h+h^2) + (1+h) + y^2+y\\
&= 2h^2 - y - 4yh + O(y^2, h^3, yh^2).
\end{align*}
Setting $y=s^2$ and seeking $h=\sigma s+O(s^2)$, the leading terms give $2\sigma^2 s^2 - s^2 = 0$, hence $\sigma^2=\tfrac12$, i.e.\ $\sigma=\pm\tfrac{1}{\sqrt{2}}$.
The correct square-root opening is confirmed by evaluating the quadratic discriminant:
\[
\frac{\partial^2\Pi}{\partial\lambda^2}\bigg|_{(1,0)} = 2\cdot\frac{\partial}{\partial\lambda}(4\lambda^3-3\lambda^2-(2y+1)\cdot 2\lambda+1)\bigg|_{(1,0)}=2(12-6-4)=4,
\]
\[
\frac{\partial\Pi}{\partial y}\bigg|_{(1,0)} = (-2\lambda^2+2y+1)\bigg|_{(1,0)}=-2+1=-1.
\]
By the Puiseux implicit function theorem for a double root (see, e.g., \cite{BrieskornKnorrer1986}), and with the same positive real square-root coordinate \(s=y^{1/2}>0\) used in Lemma~\ref{lem:local-dominance}, the two branches issuing from $(1,0)$ satisfy
\begin{equation}\label{eq:local-bifurc}
\lambda_\pm(y)=1\pm\frac{1}{\sqrt{2}}\,y^{1/2}+O(y).
\end{equation}
The coefficient is
\(\bigl(-\partial_y\Pi/(\tfrac12\partial_\lambda^2\Pi)\bigr)^{1/2}_{(1,0)}=1/\sqrt2\).
In particular, setting $y=s^2$,
\begin{equation}\label{eq:lambda-pm-local}
\lambda_+(s) = 1 + \frac{1}{\sqrt{2}}\,s + O(s^2),
\qquad
\lambda_-(s) = 1 - \frac{1}{\sqrt{2}}\,s + O(s^2).
\end{equation}
These labels are the local positive-side labels only; in the endpoint zero theorem the same analytic branches are continued to the negative endpoint scale by taking \(s=iu/m\), as stated in Theorem~\ref{thm:endpoint-zero-quantization}.

\medskip
\noindent\textbf{Spectral-radius comparison.}
For the remaining two roots near $(-1,0)$ and $(0,0)$, the implicit function theorem gives simple analytic extensions:
\begin{equation}\label{eq:other-roots-local}
\lambda_{-1}(y) = -1 + O(y),
\qquad
\lambda_0(y) = O(y).
\end{equation}
More precisely, differentiating $\Pi(\lambda,y)=\lambda^4-\lambda^3-(2y+1)\lambda^2+\lambda+y(y+1)$ with respect to $\lambda$:
\[
\partial_\lambda\Pi(\lambda,y) = 4\lambda^3-3\lambda^2-2(2y+1)\lambda+1.
\]
At $(\lambda,y)=(-1,0)$: $\partial_\lambda\Pi(-1,0) = -4-3+2+1=-4\neq 0$, so the implicit function theorem applies.
For the $y$-derivative: $\partial_y\Pi(\lambda,y)=-2\lambda^2+2y+1$, giving at $(-1,0)$: $\partial_y\Pi(-1,0)=-2+1=-1$.
Hence
\[
\frac{d\lambda_{-1}}{dy}\bigg|_{y=0} = -\frac{\partial_y\Pi(-1,0)}{\partial_\lambda\Pi(-1,0)} = -\frac{-1}{-4} = -\frac{1}{4},
\]
and
\[
\lambda_{-1}(y) = -1 - \frac{1}{4}y + O(y^2).
\]
For the root near $\lambda=0$: $\partial_\lambda\Pi(0,0)=1$ and $\partial_y\Pi(0,0)=1$, so
\[
\lambda_0(y) = -y+O(y^2).
\]
Therefore $|\lambda_{-1}(s^2)|=1+\tfrac{1}{4}s^2+O(s^4)>1$ for $s>0$ small, while $|\lambda_+(s)|=1+\tfrac{1}{\sqrt{2}}s+O(s^2)>1$, and $|\lambda_0(s^2)|=s^2+O(s^4)\to 0$ for $s\to 0^+$.

The critical comparison is $|\lambda_+(s)|$ vs $|\lambda_{-1}(s^2)|$: the bifurcating root $\lambda_+$ departs from $1$ at rate $\tfrac{1}{\sqrt{2}}s$ (first order in $s$), while the non-bifurcating root $\lambda_{-1}$ departs from $-1$ (hence $|\lambda_{-1}|$ departs from $1$) at rate $\tfrac{1}{4}s^2$ (second order in $s$):
\[
|\lambda_+(s)| - |\lambda_{-1}(s^2)| = \frac{1}{\sqrt{2}}s - \frac{1}{4}s^2 + O(s^3) > 0
\]
for $0<s<4/\sqrt{2}=2\sqrt{2}$.
For $|\lambda_-|$ and $|\lambda_0|$: $|\lambda_-(s)|=1-\tfrac{1}{\sqrt{2}}s+O(s^2)<1$ and $|\lambda_0(s^2)|=s^2+O(s^4)\to 0$.
Hence, for $s>0$ sufficiently small,
\begin{equation}\label{eq:dominance-ineqs}
\abs{\lambda_+(s)}
>
\abs{\lambda_{-1}(s^2)},
\qquad
\abs{\lambda_+(s)}
>
\abs{\lambda_-(s)},
\qquad
\abs{\lambda_+(s)}
>
\abs{\lambda_0(s^2)}.
\end{equation}
Thus the branch $\lambda_+(y^{1/2})$ is the unique maximal-modulus root for $y>0$ small, proving part~(1).

Part~(2) is immediate from
\[
\Pi(\lambda,1)=(\lambda-2)(\lambda-1)(\lambda+1)^2,
\]
since the double root has modulus \(1\) while the simple root \(\lambda=2\) has modulus \(2\).
Part~(3) is exactly Lemma~\ref{lem:ysub-exact-subdominance}: the ramification abscissa \(\alpha\) maps to \(\ySub\), the residual quadratic at \(y=\ySub\) has a non-real conjugate pair, and \eqref{eq:ysub-modulus-gap} puts that pair at strictly larger modulus than the colliding root \(\alpha\).
\end{proof}

\begin{remark}\label{rem:ysub-local-gap}
The strict inequality $\sqrt{q}>|\alpha|$ gives a positive modulus gap at \(y=y_{\mathrm{sub}}\).
By continuity of the roots of \(\Pi(\lambda,y)\) with respect to \(y\), there exists a neighborhood \(U\) of \(y_{\mathrm{sub}}\) such that, for \(y\in U\setminus\{y_{\mathrm{sub}}\}\), the two branches colliding at \(y_{\mathrm{sub}}\) remain strictly below the maximal-modulus modulus.
Thus \(y_{\mathrm{sub}}\) does not enter the maximal-modulus envelope even locally.
\end{remark}

\begin{remark}\label{rem:maximal-modulus-endpoint}
In the spectral decomposition of the polynomial sequence, the limiting zero-accumulation curve is governed by equimodular competition among maximal-modulus eigenvalues, not by the discriminant locus alone.
A branch point belongs to the maximal-modulus spectral boundary only when the colliding eigenvalue is maximal-modulus.
Theorem~\ref{thm:ep-classification} shows that $y_{\mathrm{sub}}$ fails this dominance test: the colliding eigenvalue pair at $\lambda = \alpha$ is exceeded in modulus by the complex-conjugate pair of modulus $\sqrt{q}$.
Thus $y_{\mathrm{sub}}$ is a branch point of the discriminant curve that does not enter the maximal-modulus envelope.
\end{remark}

\begin{remark}\label{rem:yang-lee-distinction}
The terminology of endpoint zeros is used here in the spectral sense of a maximal-modulus branch boundary.
The classical Yang--Lee and Fisher settings \cite{YangLee1952I,LeeYang1952II,Fisher1978,KortmanGriffiths1971} concern additional model-specific structures, which are not assumed in this article.
Thus the point $y=0$ is a square-root endpoint of the maximal-modulus boundary for this quartic family.
\end{remark}

\subsection{Real ovals and the exact real-spectral phase diagram}

The elliptic normalization also determines the complete real-root count of the quartic transfer polynomial.
The next result is a preparatory real-oval corollary; the definitive branch-value, real-root, and dominance classification is Theorem~\ref{thm:consolidated-branch-dominance-package}.

\begin{corollary}[Preparatory real-oval image diagram]\label{cor:real-phase-diagram}
Let $\beta_1 < \beta_2 < \beta_3$ be the three real roots of
\[
4\lambda^3 - 4\lambda + 1 = 0.
\]
Then the real locus of the elliptic curve
\[
E:\quad Y^2 = 4\lambda^3 - 4\lambda + 1
\]
decomposes as
\[
E(\RR)=E_{\mathrm{c}}(\RR)\sqcup E_{\mathrm{n}}(\RR),
\]
where
\[
E_{\mathrm{c}}(\RR)
=\{(\lambda,Y)\in\RR^2:\beta_1\le \lambda\le \beta_2,\ Y^2=4\lambda^3-4\lambda+1\}
\]
is a compact oval and
\[
E_{\mathrm{n}}(\RR)
=\{(\lambda,Y)\in\RR^2:\lambda\ge \beta_3,\ Y^2=4\lambda^3-4\lambda+1\}\cup\{O\}
\]
is the noncompact oval completed by the point at infinity $O$.
Under the degree-$4$ \(y\)-projection
\[
y=\lambda^2-\frac12+\frac12Y,
\]
their images are exactly
\begin{equation}\label{eq:real-oval-images}
y\bigl(E_{\mathrm{c}}(\RR)\bigr)=[y_{\mathrm{sub}},1],
\qquad
y\bigl(E_{\mathrm{n}}(\RR)\bigr)=[0,\infty).
\end{equation}
Consequently the real-root count for regular real values is obtained by applying the two-to-one covering property on the interiors of these two intervals.
The definitive table is Theorem~\ref{thm:theorem-b-real-projection-package}, and the same table is incorporated into the final branch--dominance package in Theorem~\ref{thm:consolidated-branch-dominance-package}.
\end{corollary}

\begin{proof}
The cubic $4\lambda^3-4\lambda+1$ has three distinct real roots because its derivative
\[
12\lambda^2-4=4(3\lambda^2-1)
\]
vanishes at $\lambda=\pm 1/\sqrt{3}$, with
\[
4\Bigl(-\frac{1}{\sqrt{3}}\Bigr)^3-4\Bigl(-\frac{1}{\sqrt{3}}\Bigr)+1
=1+\frac{8}{3\sqrt{3}}>0,
\]
while
\[
4\Bigl(\frac{1}{\sqrt{3}}\Bigr)^3-4\Bigl(\frac{1}{\sqrt{3}}\Bigr)+1
=1-\frac{8}{3\sqrt{3}}<0.
\]
Hence the real locus of $E$ has two connected components, one lying above the bounded interval $[\beta_1,\beta_2]$ and one lying above $[\beta_3,\infty)$ together with $O$; this is the standard topology of a real elliptic curve with three real branch points \cite{Silverman2009,Miranda1995}.

The critical points of the \(y\)-projection on $E$ are determined by $\dd y=0$, hence by
\[
(\lambda-1)(\lambda+1)(16\lambda^3-9\lambda^2+1)=0
\]
from \eqref{eq:ramification-poly}.
By Theorem~\ref{thm:ep-classification}, the cubic factor has a unique real root $\alpha\in(-1/2,0)$, and its image is $y_{\mathrm{sub}}$.
Moreover, \(4\lambda^3-4\lambda+1\) is positive at \(\lambda=-1,0,1\),
while the cubic is negative at $\lambda=1/\sqrt3$.
Hence $\beta_1<-1$, $\beta_2\in(0,1/\sqrt3)$, and $\beta_3\in(1/\sqrt3,1)$, so
\[
\beta_1<-1<\alpha<\beta_2<\beta_3<1.
\]
the real ramification points split as follows:
the compact oval contains the two critical points with $\lambda=-1$ and $\lambda=\alpha$, while the noncompact oval contains the critical point with $\lambda=1$ and the pole at infinity.

Restrict $y$ to the compact oval $E_{\mathrm{c}}(\RR)$.
This restriction is a continuous real-valued function on a compact connected one-manifold, and its only critical values are $1$ and $y_{\mathrm{sub}}$.
Therefore the image is the closed interval $[y_{\mathrm{sub}},1]$.
Similarly, on the noncompact oval the only finite critical point is $y=0$, while $y\to+\infty$ along either branch approaching $O$ because $y$ has a pole of order $4$ at infinity.
Hence the noncompact image is $[0,\infty)$.
This proves \eqref{eq:real-oval-images}.

For real $y$ away from the branch locus, each point of $y^{-1}(y)\cap E(\RR)$ is simple and contributes a distinct real root $\lambda$ of $\Pi(\lambda,y)=0$.
Moreover, each regular value in the interior of an oval image has exactly two preimages on that oval.
Thus the displayed images are precisely the component source for the four rows of Theorem~\ref{thm:theorem-b-real-projection-package}; the count itself is left there as the citable real-projection theorem.
\end{proof}

